\documentclass[12pt]{article}
\usepackage[spanish]{babel}
\usepackage{natbib}
\usepackage{url}
\usepackage[utf8x]{inputenc}
\usepackage[mediumspace,mediumqspace,squaren]{SIunits}
\usepackage{amsmath}
\usepackage{graphicx}
\graphicspath{{images/}}
\usepackage{parskip}
\usepackage{fancyhdr}
\usepackage{vmargin}
\usepackage{enumerate} % enumerados
\usepackage{diagbox}
\usepackage{booktabs}
\usepackage{float} %para usar [H]


\setmarginsrb{3 cm}{2.5 cm}{3 cm}{2.5 cm}{1 cm}{1.5 cm}{1 cm}{1.5 cm}

\title{Practico de Laboratorio N°2: Adaptación en lineas de transmisión con un stub }								% Title
\author{}								% Author
\date{2015}											% Date

\makeatletter
\let\thetitle\@title
\let\theauthor\@author
\let\thedate\@date
\makeatother

\pagestyle{fancy}
\fancyhf{}
\rhead{\theauthor}
\lhead{\thetitle}
\cfoot{\thepage}

\begin{document}

%%%%%%%%%%%%%%%%%%%%%%%%%%%%%%%%%%%%%%%%%%%%%%%%%%%%%%%%%%%%%%%%%%%%%%%%%%%%%%%%%%%%%%%%%

\begin{titlepage}
	\centering
    \vspace*{0.5 cm}
     \includegraphics[scale = 0.8]{Logo_UTN.png}\\[1.0 cm]	% University Logo
    \textsc{\LARGE Universidad Tecnológica Nacional}\\[2.0 cm]	% University Name
	\textsc{\Large Curso 3R2}\\[0.5 cm]				% Course Code
	\textsc{\large Medios de Enlace}\\[0 cm]				% Course Name
	\textrm{\large (Ing. Contreras, Luis Candelario)}\\[0.5 cm]
    \rule{\linewidth}{0.2 mm} \\[0.4 cm]
	{ \huge \bfseries \thetitle}\\
	\rule{\linewidth}{0.2 mm} \\[1 cm]
	
	\begin{minipage}{0.4\textwidth}
		\begin{flushleft} \large
			\emph{Autores:}\\
			Bianchini, Bruno\\Fichetti P. Tomás\\Guizzo, José\\Bosse, Esteban A.
			\end{flushleft}
			\end{minipage}~
			\begin{minipage}{0.4\textwidth}
			\begin{flushright} \large
			\emph{Legajo:} \\
			 63070\\62281\\62165\\62930									% Your Student Number
		\end{flushright}
	\end{minipage}\\[2 cm]
	
	{\large \thedate}\\[2 cm]
 
	\vfill
	
\end{titlepage}

%%%%%%%%%%%%%%%%%%%%%%%%%%%%%%%%%%%%%%%%%%%%%%%%%%%%%%%%%%%%%%%%%%%%%%%%%%%%%%%%%%%%%%%%%

\tableofcontents
\pagebreak

%%%%%%%%%%%%%%%%%%%%%%%%%%%%%%%%%%%%%%%%%%%%%%%%%%%%%%%%%%%%%%%%%%%%%%%%%%%%%%%%%%%%%%%%%
\section{Marco Te\'orico}
Se tiene una fuente de señal conectada a una línea de transmisión de impedancia característica $Z_O = 50 \Omega$
la que está terminada en una impedancia de carga $Z_L$.

Se desea adaptar la impedancia de carga $Z_L$ a la impedancia característica de la línea $Z_O = 50 \Omega$, con un
\textbf{stub paralelo} terminado en circuito abierto.

La frecuencia del generador es: $f = 410 MHz.$
Considerar que la velocidad de propagación en el cable coaxil es $v_P = 0.66*C$ Por lo tanto, la longitud de onda será:
$\lambda = \frac{v_P}{f}\lambda = \frac{v_P}{f}$

Las impedancias de carga $Z_L$ valen:
\begin{itemize}
 \item $Z_{L1} = 17 - j 5,3 \Omega$
 \item $Z_{L2} = 50 \Omega$
\end{itemize}

\section{Materiales Requeridos}
Los elementos utilizados en siguiente práctico son:
\begin{itemize}

 \item Generador de onda cuadrada ($f = 100 KHz; V= 3 V_{pp}$) 
 \item Vatímetro marca BIRD.
 \item Tapón de medición para vatímetro BIRD: $f = 200-500 MHz; P = 50 W$
 \item Fuente de alimentación $12 V$.
 \item Línea de transmisión $Z_O = 50 \Omega$.
 \item Conectores varios: BNC hembra-hembra, BNC triple hembra (T),UHF hembra – N macho.
 \item Impedancias de carga ($Z_{L1}$  y $Z_{L2}$).
 \item Cables varios (numerados 10, 11, 12) para interconexión (ya instalados).
\end{itemize}


\section{Procedimiento}

\subsection{Actividad N°1}
Identificar el esquema del punto 4.3 (ya está armado en el LdC).

\subsection{Actividad N°2}
Conectar la impedancia de carga $Z_{L1}$ al vatímetro.

\subsection{Actividad N°3}
Conectar el generador de señal a la frecuencia de trabajo $f = 410 MHz$.
\subsection{Actividad N°4}
Con el vatímetro, verificar la existencia de potencia $P_i$ entregada por el transmisor (debe indicar LOW) a la impedancia de carga
$Z_L$.

\subsection{Actividad N°5}
Medir la potencia incidente $P_i$  y la potencia reflejada $P_r$. Para ello se debe girar el tapón de medición del Vatímetro BIRD hacia la
carga $(P_i)$ o hacia el generador $(P_r)$, según indique la flecha en la parte superior del tapón de medición.

\textbf{En $Z_{L1}$:}

$P_i:4 $
$P_r:0.6 $

\textbf{En $Z_{L2}$:}

$P_i:3.2 $
$P_r:0.6 $

\subsection{Actividad N°6}
 Calcular el coeficiente de reflexión $\Gamma 1$  y la Relación de Onda Estacionaria sin ningún stub ($ROE1ss$) conectado a la
 línea de transmisión.

\textbf{En $Z_{L1}$:}

$\Gamma 1_{ss}:0.194$
$ROE1ss: 1.48$

\textbf{En $Z_{L2}$:}

$\Gamma 1_{ss}:0.2421$

$ROE1ss:1.64$

\subsection{Actividad N°7}
Conectar el adaptador BNC tipo T luego del vatímetro.

\subsection{Actividad N°8}
Conectar la impedancia de carga mediante el segmento de línea calculado y en el extremo libre, conectar el stub
adecuado de longitud $L$.

\subsection{Actividad N°9}
Medir la potencia incidente $P_i$ y la potencia reflejada $P_r$.

\textbf{En $Z_{L1}$:}

$P_i:4.4 $
$P_r:0.1 $

\textbf{En $Z_{L2}$:}

$P_i:3.2 $
$P_r:0 $

\subsection{Actividad N°10}
 Calcular el coeficiente de reflexión $\Gamma 1$  y la Relación de Onda Estacionaria con el stub
 ($ROE1cs$) conectado a la línea de transmisión.
\textbf{En $Z_{L1}$:}

 $\Gamma 1_{cs}:0$
 $ROE1ss: 1$

 \textbf{En $Z_{L2}$:}
 
 $\Gamma 2_{cs}:0$
 $ROE2ss: 1$
\subsection{Actividad N°11}
 Repetir desde el punto 2) para la carga $Z_{L2}$.
 
 Est\'an resueltos al final de cada actividad.
 
\section{Cuestionario}
\begin{itemize}

\item ¿Qué implica que la línea de transmisión esté adaptada?

Implica que no va a existir onda reflejada en la linea de transmisión, esto hace que la carga pueda aprovechar toda la potencia del generador.
\item ¿Qué resultado se obtendría si el stub para $Z_{L1}$  fuera más largo?

Existir\'ia una onda estacionaria reflejada, dejando la linea de transmisión no adaptada y produce un movimiento en sentido horario en el \'abaco de smith en el largo del stub.
\item ¿Qué resultado se obtendría si el stub para $Z_{L1}$  fuera más corto?

Produce el mismo fen\'omeno que la linea mas corta pero esta hace que nos movamos en el \'abaco en sentido anti-horario.
\item ¿Por qué para $ZL2$ resulta $\Gamma 2 = 0$ y $ROE2 = 1$?

Esto sucede porque la linea se encuentra perfectamente adaptada, no habiendo onda reflejada y por ende siendo el $ROE2 = 1$.
\end{itemize}

\end{document}